% 1. VALORES
\section{Valores Fundamentais}

Os valores do Clube de Matemática orientam nossas interações, decisões e a cultura do grupo.
Eles refletem nosso compromisso com um ambiente de aprendizado inclusivo, seguro e eficaz,
onde todo cadete pode crescer e contribuir.

\subsection{Apoio Mútuo e Responsabilidade Compartilhada}

Cada membro contribui para o esforço coletivo e reconhece seu papel como essencial para o
sucesso do grupo. Isso inclui não apenas comprometimento individual, mas também a disposição
de apoiar e encorajar colegas que estejam enfrentando dificuldades seja com um conceito
matemático ou com a permanência na 42.

\subsection{Comunicação Aberta e Construtiva}

Criar um espaço seguro onde os cadetes possam tirar dúvidas sem julgamento é central para
o clube. Ninguém deve sentir vergonha por não entender algo --- a matemática é uma jornada
e toda pergunta é válida. O feedback entre membros deve ser honesto, respeitoso e voltado
ao crescimento.

\subsection{Respeito às Diferenças}

Os cadetes da 42 vêm de formações muito diversas: alguns nunca viram Cálculo, outros têm
graduação em exatas. Essa diversidade é uma força. O clube valoriza diferentes formas de
pensar e aprender, criando um ambiente onde todos se sintam confortáveis para contribuir
independentemente do seu ponto de partida.

\subsection{Pertencimento e Presença}

Um dos objetivos centrais do clube é fazer com que o cadete sinta que tem um lugar na
42 --- pessoas que esperam por ele, um grupo ao qual pertence. Presença regular nas sessões
não é apenas uma questão de aprendizado, mas de fortalecer os laços que mantêm a
comunidade unida.

\subsection{Curiosidade e Persistência}

A matemática exige tempo e prática. O clube celebra o processo, não apenas o resultado.
Errar faz parte do aprendizado e encarar desafios difíceis com persistência é mais valioso
do que chegar rápido à resposta certa.
